\chapter{\label{III_2}Adapter ces fonctionnalités à tous les clients (passer du cas d’étude à la généralisation} 

L'éditeur logiciel aurait un rôle à jouer dans la mise à disposition de nouvelles fonctionnalités innovantes aux musées et centres d'archives, d'autant plus ceux qui préservent des archives audiovisuelles. En effet, "les institutions de mémoire contenant des contenus audiovisuels sont prédestinées à explorer les avantages du ML ou de l'intelligence hybride pour la mise en valeur et l'accès à leurs contenus audiovisuels. [...] dans ce cas, la collaboration avec des partenaires, tels que des prestataires de services ou des portails, s'impose.” Nous avons vu comment proposer un pipeline automatique à partir d'un exemple de corpus audiovisuel bien précis. Toutefois, ce pipeline doit pouvoir être adapté à tous types de corpus et à toutes institutions hébérgées par le logiciel.  




\section{\label{III_2_A}Logique produit vs logique patrimoniale}

La proposition d'une nouvelle fonctionnalité au sein d'un logiciel de gestion des collections impose de concilier deux logiques distinctes : la logique industrielle de l'éditeur logiciel et la logique de préservation et de gestion des institutions patrimoniales. Ainsi, l'objectif pour l'éiteur n'est pas seulement de proposer des fonctionnalités performantes, mais surtout de proposer des fonctionnalités adaptatives à tous les types de collections potentiellement gérées par les isntitutions clientes. 
Le premier obstacle à l'adaptation d'une fonctionnalité fondée sur l'IA comme c'est le cas de l'indexation automatisée réside dans le paramétrage des modèles de visions par ordinateur. Comme nous l'avons évoqué dans notre étude, un pipeline de segmentation (pyscenedetect) et de détection d'objets (Yolov8) nécessite des ajustements fins en fonction du corpus étudié, notamment en ce qui concerne le seuil de confiance ou le seuil de séquencage. La généralisation de ce pipeline à une clientèle large suppose de renoncer aux modèles lourds que nous avons évoqué tels que les visions transformers, fondés sur le deep learning, qui nécessitent des paramétrages locaux lourds. La tendance s'oriente vers des architectures légères mais hautement interactives afin de permettre à l'utilisateur d'être intégré au fonctionnement de l'algorithme (référence?). Objectif : pouvoir garder le contrôle sur le paramétrage.
Le choix de ce type d'architecture pose la question de l'hébergement des données, pour pouvoir garder le contrôle sur celles-ci. En effet, l'éditeur de logiciel de gestion des collections est reponsable des données clients. Skinsoft propose dans la plupart des cas d'héberger les données de ses clients, afin d'éviter d'installer un serveur de données chez chaque client individuellement et pouvoir facilement gérer les problèmes potentiels. La plupart des outils automatiques disponibles et prêt à l'emploi sont hébergés sur des nuages externes (\textit{clouds}). Dans le cas de l'adoption d'un tel outil en nuage, cela suppose que les données clients sont "envoyées" sur un serveur externe appartenant à une grande entreprise, sur lequel l'éditeur de logiciel n'aurait qu'un contrôle limité, créant ainsi une "dépendance". \footcite{zotero-item-584} En tant qu'éditeur, le rôle de Skinsoft, par exemple, serait de s'approprier un algorithme déjà existant et de le développer pour pouvoir l'adapter à ses clients. De fait, les données seraient entièrement hébérgées par l'éditeur permettant ainsi de gérer les potentiels problèmes sur les données entrantes au sein des algorithmes et les données sortantes. 

De plus, au delà des contraintes posées par de nouvelles fonctionnalités automatisées, se pose la question de la pertinence du déploiement de tels outils, en lien avec les besoins profesionnels actuels. En effet, l'éditeur de logiciel est animé par une logique commerciale qui devance les pratiques ancestrales de gestion des collections. La posture de l'éditeur, en prise avec l'innovation pour devancer toute concurrence contraste donc avec les logiques institutionnelles et patrimoniales de ses clients. Dans son guide pour une IA responsable dans le secteur culturel, le Ministère de la Culture rappelle que dans le cadre d'une mise en place de fonctionnalités fondées sur l'IA, il s'agit "d’examiner le besoin à couvrir et de vérifier la capacité des systèmes d’IA à satisfaire ce besoin dans des conditions optimales" ; le cas échéant, de favoriser des alternatives moins consommatrices que l’IA si elles sont suffisantes à répondre au besoin. » \footcite{2025a}
Dans le cas de Skinsoft, si certains clients sont demandeurs d'une automatisation de certaines tâches réalisables au sien d'un système de gestion des collections, d'autres les refusent catégoriquement de peur de perdre le contrôle sur les données existantes et celles qui seront produites. Le pipeline proposé dans notre étude, fondé sur la computer vision est une alternative intéressante puisqu'elle permet de configurer pleinement les tâches automatisées et de garder le contrôle sur ce qui est produit. Les risques associés habituellement aux modèles d'IA tels que les risque d'hallucination, d'erreurs, de biais, seraient amoindris dans le cadre de ce pipeline, puisque les outils seraient intégrés à un logiciel professionnels comme facilitateurs et propositions, mais jamais comme décisionnaires. 
exemple du NLP : nous avons proposé le NLP comme un module que l'utilisateur peut désactiver s'il le souhaite. 
Les professionnels pourraient mettre à profit l'IA pour un double objectif : la préservation et l'accessibilité des collections. \footcite{zotero-item-640}. C'est exactement le cas pour l'indexation automatique.



Deux idées : 
- adapter via des pipelines plus lourds mais contraintes d'infrastructures (Clariah Dane) le pipeline que nous avons proposé à l'aide de la computer vision nécessite un paramétrage pour chaque corpus pour le seuil de séquencage. et yolo ? comment l'adapter ? 


Proposer des outils plus performants que ceux proposés dans cette étude afin de pouvoir s’adapter à tous les types de corpus, tout en prenant en compte les contraintes d’infrastructures. Mais : contraintes d’infrastructures. Skinsoft héberge les données clients donc amener un algorithme lourd en local.  

Transformers : Revisiting Human-in-the-Loop Object Retrieval with Pre-Trained Vision Transformers  K. Zaher, O. Buisson et A. Joly ICPR 2026 - 28th International Conference on Pattern Recognition, 2026le 


En 2025, le ministère de la culture identifie "l'enrichissement de la connaissances et des savoirs" comme un usage potentiel de l'IA en secteur culturel avec par exemple "de l’aide au catalogage, à l’indexation automatique des collections, au tri et au classement des documents et données destinés à être archivés, à la transcription des textes manuscrits, à la reconnaissance de formes, à la diffusion et à la valorisation des contenus culturels et patrimoniaux quel que soit leur format (image, son…)". \footcite{2025a}
\section{\label{III_2_B}Jeux de données patrimoniaux} 

The ‘collections as data’ movement seeks to make digital cultural heritage materials available for computational use, encouraging LAM institutions to provide their holdings in machine-actionable formats.Footnote25 For archives, this shift can be valuable, enabling new forms of analysis and engagement. However, as Caswell reminds us, ‘the archive’ in abstract computational discourse is not the same as ‘archives’ as institutions, records, and practices.Footnote26 Treating archives exclusively as data risks flattening their complex provenance, legal status, and social meanings. In this article we therefore adopt a more cautious framing of collections as/for data: we acknowledge that digital surrogates and descriptive metadata can be structured for computational purposes, while insisting that such structuring must remain accountable to archival principles and to the communities whose lives and histories are represented in the records. \footcite{colavizza2026}

La validité d'un modèle d'apprentissage automatique repose sur la qualité et la pertinence des données d'entraînement. Si notre étude a montré la potentielle efficacité d'un pipeline de segmentation puis d'indexation sur un corpus d'images animées spécifique, la généralisation de cette fonctionnalité à un nombre large d'institutions se heurte à l'hétérogénéité des fonds audiovisuels patrimoniaux. 
C'est essentiellement la reconnaissance d'objets fondé sur YoLov8 qui nécessite un jeu de données d'entraînement. L'outil est actuellement entraîné sur des jeux de données d'images génériques, différentes d'images patrimoniales. Les collections conservées au sein des musées et des centres d'archives audiovisuelles se caractérisent en effet par des spécificités visuelles et sémantiques uniques : supports anciens dégradés, esthétiques cinématographiques variées, vocabulaire documentaire hautement spécialisé. Un modèle d'IA entraîné sur un jeu de données restreint risque de produire un surapprentissage (\textit{overfitting}) peu adaptable à la variété des collections gérées par les différents clients de l'éditeur. Pour garantir la robustesse des fonctionnalités automatiques, il s'avère donc indispensable de s'appuyer sur des jeux de données d'entraînement davantage représentatif.
Pour constituer ces corpus d'entraînement à grande échelle sans dépendre des données privées des institutions, la stratégie la plus pertinente réside dans la réutilisation des ressources publiques sous licence libre. 
Dans cette optique, l'entraînement et l'affinage (\textit{fine-tuning}) des modèles de vision par ordinateur peuvent s'appuyer sur d'immenses réservoirs de données ouvertes. Le projet Dataset, créé par l'INA, met notamment à disposition des corpus de données audiovisuelles  "destiné à la mise au point, l’expérimentation et l’évaluation d’outils de recherche et d’analyse de contenus multimédias " \footcite{zotero-item-754}. Ces corpus contiennent des images animées issues de la télévision : journaux télévisés, actualités, captations... Or, ce jeu de données n'entraîne pas nécessairement le modèle sur des images dégradées, en pellicule, en noir et blanc. 
Les Prelinger archives, par exemple, réunissent des archives audiovisuelles américaines de tous types, dénuées de droits d'auteur, y compris des films amateurs ou éducatifs par exemple. \footcite{zotero-item-756}
wikimedia commons 
L'exploitation des API et des corpus sous licence ouverte d'intégrateurs majeurs tels qu'Europeana, le Portail Open Data de la Plateforme Ouverte du Patrimoine (POP) du Ministère de la Culture, ou les fonds numérisés de Gallica (BnF).

Double entraînement : données + règles métier (FRBR) : pour le NLP : LLM

L'entraînement des modèles reposant sur un fonds unique d'un client donné, comme nous l'avons exposé en seconde partie, poserait plusieurs difficultés pour la généralisation de la fonctionnalité. D'une part, les données d'une seule institution ne sont pas suffisantes ; d'autre part, cela exige un accord contractuel explicite de cession de droits d'usage sur les données. La création d'un modèle mutuel suppose une gouvernance transparente, où chaque institution accepte d'alimenter un réservoir partagé de \textit{Ground Truth} dans une logique de mutualisation des ressources.

Enfin, cette réflexion sur les jeux de données d'entraînement résonne directement avec les autres briques technologiques en cours de spécification au sein de l'entreprise, notamment la recherche en langage naturel et la recherche visuelle par similarité. La construction d'un jeu de données patrimonial universel — associant visuels sous licence libre, descripteurs FRBR et thésaurus normés — profite à l'ensemble du système d'information :
\begin{enumerate}
	\item Les métadonnées textuelles ouvertes renforcent le \textit{mapping} sémantique de la recherche en langage naturel.
	\item Les paires images-concepts réutilisables affinent les réseaux de neurones dédiés à la recherche par l'image et à la segmentation vidéo.
\end{enumerate}

En plaçant la réutilisation des données ouvertes et le principe des \textit{Collections as Data} au cœur de sa R\&D, Skinsoft garantit la soutenabilité technique de ses fonctionnalités automatiques tout en respectant scrupuleusement le cadre légal et éthique de ses clients patrimoniaux.



choix d'appui sur des données produites directement par des institutions culturelles. utiliser exclusivement des sources ouvertes. pour les images : entrainement sur des images disponibles sous licence ouverte : european, POP + textes de lois, méthodologie des gestion des collections, thesaurus, normes de catalogage...R
Il faut l'accord client et les données d'un seul client ne sont pas suffisantes.
\section{\label{III_2_B}Jeux de données patrimoniaux} 

The ‘collections as data’ movement seeks to make digital cultural heritage materials available for computational use, encouraging LAM institutions to provide their holdings in machine-actionable formats.Footnote25 For archives, this shift can be valuable, enabling new forms of analysis and engagement. However, as Caswell reminds us, ‘the archive’ in abstract computational discourse is not the same as ‘archives’ as institutions, records, and practices.Footnote26 Treating archives exclusively as data risks flattening their complex provenance, legal status, and social meanings. In this article we therefore adopt a more cautious framing of collections as/for data: we acknowledge that digital surrogates and descriptive metadata can be structured for computational purposes, while insisting that such structuring must remain accountable to archival principles and to the communities whose lives and histories are represented in the records. \footcite{colavizza2026}

La validité d'un modèle d'apprentissage automatique repose sur la qualité et la pertinence des données d'entraînement. Si notre étude a montré la potentielle efficacité d'un pipeline de segmentation puis d'indexation sur un corpus d'images animées spécifique, la généralisation de cette fonctionnalité à un nombre large d'institutions se heurte à l'hétérogénéité des fonds audiovisuels patrimoniaux. 
C'est essentiellement la reconnaissance d'objets fondé sur YoLov8 qui nécessite un jeu de données d'entraînement. L'outil est actuellement entraîné sur des jeux de données d'images génériques, différentes d'images patrimoniales. Les collections conservées au sein des musées et des centres d'archives audiovisuelles se caractérisent en effet par des spécificités visuelles et sémantiques uniques : supports anciens dégradés, esthétiques cinématographiques variées, vocabulaire documentaire hautement spécialisé. Un modèle d'IA entraîné sur un jeu de données restreint risque de produire un surapprentissage (\textit{overfitting}) peu adaptable à la variété des collections gérées par les différents clients de l'éditeur. Pour garantir la robustesse des fonctionnalités automatiques, il s'avère donc indispensable de s'appuyer sur des jeux de données d'entraînement davantage représentatif.
Pour constituer ces corpus d'entraînement à grande échelle sans dépendre des données privées des institutions, la stratégie la plus pertinente réside dans la réutilisation des ressources publiques sous licence libre. 
Dans cette optique, l'entraînement et l'affinage (\textit{fine-tuning}) des modèles de vision par ordinateur peuvent s'appuyer sur d'immenses réservoirs de données ouvertes. Le projet Dataset, créé par l'INA, met notamment à disposition des corpus de données audiovisuelles  "destiné à la mise au point, l’expérimentation et l’évaluation d’outils de recherche et d’analyse de contenus multimédias " \footcite{zotero-item-754}. Ces corpus contiennent des images animées issues de la télévision : journaux télévisés, actualités, captations... Or, ce jeu de données n'entraîne pas nécessairement le modèle sur des images dégradées, en pellicule, en noir et blanc. 
Les Prelinger archives, par exemple, réunissent des archives audiovisuelles américaines de tous types, dénuées de droits d'auteur, y compris des films amateurs ou éducatifs par exemple. \footcite{zotero-item-756}
wikimedia commons 
L'exploitation des API et des corpus sous licence ouverte d'intégrateurs majeurs tels qu'Europeana, le Portail Open Data de la Plateforme Ouverte du Patrimoine (POP) du Ministère de la Culture, ou les fonds numérisés de Gallica (BnF).

Double entraînement : données + règles métier (FRBR) : pour le NLP : LLM

L'entraînement des modèles reposant sur un fonds unique d'un client donné, comme nous l'avons exposé en seconde partie, poserait plusieurs difficultés pour la généralisation de la fonctionnalité. D'une part, les données d'une seule institution ne sont pas suffisantes ; d'autre part, cela exige un accord contractuel explicite de cession de droits d'usage sur les données. La création d'un modèle mutuel suppose une gouvernance transparente, où chaque institution accepte d'alimenter un réservoir partagé de \textit{Ground Truth} dans une logique de mutualisation des ressources.

Enfin, cette réflexion sur les jeux de données d'entraînement résonne directement avec les autres briques technologiques en cours de spécification au sein de l'entreprise, notamment la recherche en langage naturel et la recherche visuelle par similarité. La construction d'un jeu de données patrimonial universel — associant visuels sous licence libre, descripteurs FRBR et thésaurus normés — profite à l'ensemble du système d'information :
\begin{enumerate}
	\item Les métadonnées textuelles ouvertes renforcent le \textit{mapping} sémantique de la recherche en langage naturel.
	\item Les paires images-concepts réutilisables affinent les réseaux de neurones dédiés à la recherche par l'image et à la segmentation vidéo.
\end{enumerate}

En plaçant la réutilisation des données ouvertes et le principe des \textit{Collections as Data} au cœur de sa R\&D, Skinsoft garantit la soutenabilité technique de ses fonctionnalités automatiques tout en respectant scrupuleusement le cadre légal et éthique de ses clients patrimoniaux.



choix d'appui sur des données produites directement par des institutions culturelles. utiliser exclusivement des sources ouvertes. pour les images : entrainement sur des images disponibles sous licence ouverte : european, POP + textes de lois, méthodologie des gestion des collections, thesaurus, normes de catalogage...R
Il faut l'accord client et les données d'un seul client ne sont pas suffisantes.


